\documentclass[12pt]{article}

\input{~/School/Notes/header.tex}

\title{Домашнее задание по математическому анализу}
\author{Денисов Илья}
\date{}

\begin{document}
	\maketitle

	\section*{Задача 2}

	\noindent
	\textbf{Условие.} Решить уравнение $(2x+1)y' = 4x + 2y$.

	\noindent
	\textbf{Решение.} Приводим к линейному виду ($x \ne -\tfrac{1}{2}$):
	\[
		y' - \frac{2y}{2x+1} = \frac{4x}{2x+1}.
	\]
	Полагаем $y = uv$, $y' = u'v + uv'$:
	\[
		u'v + u\left( v' - \frac{2v}{2x+1} \right) = \frac{4x}{2x+1}.
	\]
	Берём частное решение уравнения в скобках:
	\[
		v' - \frac{2v}{2x+1} = 0 \quad \Longrightarrow \quad \frac{dv}{v} = \frac{2\,dx}{2x+1} \quad \Longrightarrow \quad v = 2x+1.
	\]
	Подставляем во второе уравнение:
	\[
		u'(2x+1) = \frac{4x}{2x+1} \quad \Longrightarrow \quad u' = \frac{4x}{(2x+1)^2}.
	\]
	Делаем замену $t = 2x+1$, $x = \dfrac{t-1}{2}$:
	\[
		\int \frac{4x\,dx}{(2x+1)^2} = \int \frac{t-1}{t^2}\,dt = \ln|t| + \frac{1}{t} + C = \ln|2x+1| + \frac{1}{2x+1} + C.
	\]
	Значит $u = \ln|2x+1| + \dfrac{1}{2x+1} + C$, и
	\[
		y = uv = (2x+1)\left( \ln|2x+1| + C \right) + 1.
	\]

	\medskip
	\noindent
	\textbf{Ответ.}
	\[
		\boxed{y = (2x+1)\bigl(\ln|2x+1| + C\bigr) + 1}
	\]

	%==================================================
	\section*{Задача 3}

	\noindent
	\textbf{Условие.} Решить уравнение $x^2y' + xy + 1 = 0$.

	\noindent
	\textbf{Решение.} Приводим к линейному виду ($x \ne 0$):
	\[
		y' + \frac{y}{x} = -\frac{1}{x^2}.
	\]
	Ищем $y = uv$:
	\[
		u'v + u\left( v' + \frac{v}{x} \right) = -\frac{1}{x^2}.
	\]
	Частное решение уравнения в скобках:
	\[
		v' + \frac{v}{x} = 0 \quad \Longrightarrow \quad \frac{dv}{v} = -\frac{dx}{x} \quad \Longrightarrow \quad v = \frac{1}{x}.
	\]
	Тогда
	\[
		\frac{u'}{x} = -\frac{1}{x^2} \quad \Longrightarrow \quad u' = -\frac{1}{x} \quad \Longrightarrow \quad u = -\ln|x| + C.
	\]
	Итого
	\[
		y = uv = \frac{C - \ln|x|}{x}.
	\]

	\medskip
	\noindent
	\textbf{Ответ.}
	\[
		\boxed{y = \dfrac{C - \ln|x|}{x}}
	\]

	%==================================================
	\section*{Задача 4}

	\noindent
	\textbf{Условие.} Решить задачу Коши $xy^2y' = x^2 + y^3$, $y(1) = 0$.

	\noindent
	\textbf{Решение.} Уравнение нелинейно относительно $y$, но становится линейным после замены $t = y^3$, так как $t' = 3y^2y'$:
	\[
		\frac{x}{3}t' = x^2 + t \quad \Longrightarrow \quad t' - \frac{3}{x}t = 3x.
	\]
	Решаем это линейное уравнение подстановкой $t = pq$, $t' = p'q + pq'$:
	\[
		p'q + p\left( q' - \frac{3q}{x} \right) = 3x.
	\]
	Частное решение уравнения в скобках:
	\[
		q' - \frac{3q}{x} = 0 \quad \Longrightarrow \quad \frac{dq}{q} = \frac{3\,dx}{x} \quad \Longrightarrow \quad q = x^3.
	\]
	Подставляем:
	\[
		p'x^3 = 3x \quad \Longrightarrow \quad p' = \frac{3}{x^2} \quad \Longrightarrow \quad p = -\frac{3}{x} + C.
	\]
	Значит
	\[
		t = pq = \left( C - \frac{3}{x} \right)x^3 = Cx^3 - 3x^2,
	\]
	и, возвращаясь к исходной функции,
	\[
		y^3 = Cx^3 - 3x^2.
	\]
	Из условия $y(1) = 0$: $C - 3 = 0 \Rightarrow C = 3$.

	\medskip
	\noindent
	\textbf{Ответ.}
	\[
		\boxed{y^3 = 3x^2(x-1)}
	\]

	%==================================================
	% \section*{Задача 5}

	% \noindent
	% \textbf{Условие.} Решить уравнение $x(e^y - y') = 2$.

	% \noindent
	% \textbf{Решение.} Перепишем уравнение как $xe^y - xy' = 2$. Оно нелинейно относительно $y$, но становится линейным после замены $z = e^{-y}$: тогда $z' = -e^{-y}y'$, откуда $y' = -\dfrac{z'}{z}$, $e^y = \dfrac{1}{z}$. Подставляя,
	% \[
	% 	\frac{x}{z} + \frac{xz'}{z} = 2 \quad \Longrightarrow \quad xz' - 2z = -x \quad \Longrightarrow \quad z' - \frac{2z}{x} = -1.
	% \]
	% Решаем подстановкой $z = pq$, $z' = p'q + pq'$:
	% \[
	% 	p'q + p\left( q' - \frac{2q}{x} \right) = -1.
	% \]
	% Частное решение уравнения в скобках:
	% \[
	% 	q' - \frac{2q}{x} = 0 \quad \Longrightarrow \quad \frac{dq}{q} = \frac{2\,dx}{x} \quad \Longrightarrow \quad q = x^2.
	% \]
	% Подставляем:
	% \[
	% 	p'x^2 = -1 \quad \Longrightarrow \quad p' = -\frac{1}{x^2} \quad \Longrightarrow \quad p = \frac{1}{x} + C.
	% \]
	% Значит
	% \[
	% 	z = pq = \left( C + \frac{1}{x} \right)x^2 = Cx^2 + x,
	% \]
	% и, возвращаясь к $y$,
	% \[
	% 	e^{-y} = x + Cx^2.
	% \]

	% \medskip
	% \noindent
	% \textbf{Ответ.}
	% \[
	% 	\boxed{y = -\ln\bigl|x + Cx^2\bigr|}
	% \]

	% %==================================================
	% \section*{Задача 6}

	% \noindent
	% \textbf{Условие.} Решить уравнение $y(x) = \displaystyle\int_0^x y(t)\,dt + x + 1$.

	% \noindent
	% \textbf{Решение.} По теореме Барроу продифференцируем обе части по $x$:
	% \[
	% 	y'(x) = y(x) + 1.
	% \]
	% Начальное условие получаем, подставив $x = 0$ в исходное уравнение: $y(0) = 0 + 0 + 1 = 1$.

	% Решаем линейное уравнение $y' - y = 1$ подстановкой $y = uv$, $y' = u'v + uv'$:
	% \[
	% 	u'v + u(v' - v) = 1.
	% \]
	% Частное решение уравнения в скобках:
	% \[
	% 	v' - v = 0 \quad \Longrightarrow \quad v = e^x.
	% \]
	% Подставляем:
	% \[
	% 	u'e^x = 1 \quad \Longrightarrow \quad u' = e^{-x} \quad \Longrightarrow \quad u = -e^{-x} + C.
	% \]
	% Итого
	% \[
	% 	y = uv = \left( C - e^{-x} \right)e^x = Ce^x - 1.
	% \]
	% Из условия $y(0) = 1$: $C - 1 = 1 \Rightarrow C = 2$.

	% \medskip
	% \noindent
	% \textbf{Ответ.}
	% \[
	% 	\boxed{y = 2e^x - 1}
	% \]
\end{document}